\section{Основная лемма о линейной зависимости}

\begin{lemma}[Основная лемма о линейной зависимости]
Пусть $S_1 = \{a_1, a_2, \dots, a_k\}$ и $S_2 = \{b_1, b_2, \dots, b_m\}$ --- две системы векторов из линейного пространства $V$. Если система $S_1$ линейно независима (ЛНЗ) и каждый вектор из $S_1$ линейно выражается через векторы из $S_2$, то количество векторов в первой системе не превосходит количества векторов во второй, то есть $k \le m$.
\end{lemma}

\begin{proof}
Докажем утверждение от противного. Предположим, что $k > m$.
По условию, каждый вектор $a_j$ ($j = 1, \dots, k$) линейно выражается через векторы системы $S_2$:
\[
a_j = \alpha_{1j}b_1 + \alpha_{2j}b_2 + \dots + \alpha_{mj}b_m = \sum_{i=1}^m \alpha_{ij}b_i.
\]
Рассмотрим линейную комбинацию векторов системы $S_1$, равную нулевому вектору:
\[
x_1 a_1 + x_2 a_2 + \dots + x_k a_k = 0.
\]
Подставим выражения для $a_j$ через $b_i$:
\[
\sum_{j=1}^k x_j \left( \sum_{i=1}^m \alpha_{ij}b_i \right) = 0 \quad \Rightarrow \quad \sum_{i=1}^m \left( \sum_{j=1}^k \alpha_{ij}x_j \right) b_i = 0.
\]
Чтобы эта комбинация была верной, достаточно потребовать, чтобы коэффициенты при всех $b_i$ были равны нулю. Это приводит к однородной системе линейных алгебраических уравнений (СЛАУ) относительно неизвестных $x_1, \dots, x_k$:
\[
\begin{cases}
\alpha_{11}x_1 + \alpha_{12}x_2 + \dots + \alpha_{1k}x_k = 0 \\
\alpha_{21}x_1 + \alpha_{22}x_2 + \dots + \alpha_{2k}x_k = 0 \\
\vdots \\
\alpha_{m1}x_1 + \alpha_{m2}x_2 + \dots + \alpha_{mk}x_k = 0
\end{cases}
\]
Эта система содержит $m$ уравнений и $k$ неизвестных. По нашему предположению, $k > m$ (число неизвестных строго больше числа уравнений). Известно, что такая однородная система всегда имеет нетривиальное решение (то есть существует набор $x_1, \dots, x_k$, не все равные нулю, обращающий систему в верное равенство).

Это означает, что существует нетривиальная линейная комбинация векторов $a_1, \dots, a_k$, равная нулю, что противоречит условию леммы о том, что система $S_1$ линейно независима. 
Следовательно, наше предположение $k > m$ неверно, и должно выполняться $k \le m$.
\hfill$\blacksquare$
\end{proof}
